\chapter{Perkalian: Langkah Pertama}\label{ch:g2:mult}

Empat kotak berisi tiga kue: alih-alih $3 + 3 + 3 + 3$, kita menulis
$4 \times 3$. Perkalian adalah jalan pintas untuk menjumlahkan
bilangan yang sama berkali-kali --- dan gambar terbaiknya adalah
susunan baris dan kolom yang rapi.

\section{Menjumlahkan bilangan yang sama}

\begin{definition}[Perkalian]\label{def:g2:mult:def}
$4 \times 3$ berarti ``empat kali tiga'':
\[
4 \times 3 = 3 + 3 + 3 + 3 = 12 .
\]
Tandanya $\times$, dibaca ``kali''. Hasilnya disebut \emph{hasil
kali}\index{hasil kali} (perkalian tumbuh dewasa pada \cref{ch:g3:mult}).
\end{definition}

\begin{omfigure}
\begin{tikzpicture}[scale=0.5]
  \foreach \r in {0,1,2,3}
    \foreach \c in {0,1,2}
      \fill[omDef] (\c,-\r) circle (0.28);
  \node[right, font=\small] at (3.0,-1.5)
    {$4$ baris berisi $3$ titik: $4 \times 3 = 12$};
\end{tikzpicture}

{\small Sebuah \emph{susunan}: $4$ baris berisi $3$. Membilang per
baris memberi $4 \times 3$; titik yang sama dibilang per kolom
memberi $3 \times 4$ --- jadi urutannya tidak berpengaruh.}
\end{omfigure}

\begin{example}[Urutan tidak berpengaruh]\label{ex:g2:mult:order}
$4 \times 3$ dan $3 \times 4$ membilang $12$ titik yang sama, satu
per baris, satu per kolom. Memiringkan susunannya membuktikan hal
itu. Jadi menghafal satu \omterm{def:g2:mult:def}{hasil kali} memberikan kembarannya cuma-cuma.
\end{example}

\section{Membilang loncat dan tabel yang mudah}

\begin{method}[Menyusun tabel dengan membilang loncat]\label{met:g2:mult:skip}
Untuk mencari $5 \times 4$, bilanglah empat-empat sebanyak lima kali:
$4, 8, 12, 16, 20$. Membilang loncat di garis bilangan \emph{adalah}
tabelnya.
\begin{itemize}
  \item Tabel $2$: $2, 4, 6, 8, 10, \dots$ (\omterm{def:g1:addition:def}{jumlah} kembar);
  \item tabel $5$: $5, 10, 15, 20, 25, \dots$ (berakhir dengan $5$
        atau $0$);
  \item tabel $10$: $10, 20, 30, \dots$ (cukup tambahkan sebuah \omterm{def:g1:counting:zero}{nol}).
\end{itemize}
\end{method}

\begin{omfigure}
\begin{tikzpicture}[scale=0.6]
  \draw[-Stealth] (-0.4,0) -- (11.4,0);
  \foreach \i in {0,...,11} \draw (\i,-0.1) -- (\i,0.1);
  \foreach \i in {0,2,4,6,8,10}
    \node[below, font=\footnotesize] at (\i,-0.12) {\i};
  \foreach \i in {0,1,2,3,4}
    \draw[-Stealth, omDef, thick] (2*\i,0.28) .. controls
      (2*\i+1,0.85) .. (2*\i+2,0.28);
  \node[omDef, font=\small] at (5.5,1.1) {lompatan $2$};
\end{tikzpicture}

{\small Tabel $2$ sebagai lompatan di garis bilangan: $2, 4, 6, 8, 10$.
Lima lompatan $2$ mencapai $10$, jadi $5 \times 2 = 10$.}
\end{omfigure}

\begin{example}[Kembar dan separuh]\label{ex:g2:mult:double}
Mengalikan dengan $2$ berarti melipatduakan: $2 \times 7 = 14$.
Kebalikannya, membagi dua, membatalkannya: separuh dari $14$ adalah
$7$. \omterm{def:g1:addition:def}{Jumlah} kembar yang dihafal membuat seluruh tabel $2$ menjadi
seketika --- dan menyiapkan pembagian (\cref{ch:g3:division}).
\end{example}

\section{Perkalian dalam soal cerita}

\begin{example}\label{ex:g2:mult:problem}
``Satu kotak berisi $5$ telur. Berapa telur di dalam $4$ kotak?''
Empat kelompok berisi lima: $4 \times 5 = 20$ (bilang loncat:
$5, 10, 15, 20$). Ada $20$ telur.
\end{example}

\section{Latihan}

\begin{exercise}[$\star$]\label{exo:g2:mult:1}
Tulislah sebagai perkalian, lalu hitunglah: $2 + 2 + 2$; \quad
$5 + 5 + 5 + 5$; \quad $10 + 10 + 10$.
\end{exercise}

\begin{exercise}[$\star$]\label{exo:g2:mult:2}
Gambarlah susunan $3$ baris berisi $5$ titik. Perkalian apa yang
ditunjukkannya? Dan kalau dimiringkan?
\end{exercise}

\begin{exercise}[$\star$]\label{exo:g2:mult:3}
Bilanglah loncat untuk mengisi tabel $2$ sampai $2 \times 6$; tabel
$5$ sampai $5 \times 6$.
\end{exercise}

\begin{exercise}[$\star$]\label{exo:g2:mult:4}
Hitunglah: $2 \times 7$; \quad $5 \times 3$; \quad $10 \times 4$;
\quad $2 \times 9$; \quad $5 \times 6$.
\end{exercise}

\begin{exercise}[$\star$]\label{exo:g2:mult:5}
Hitunglah kembarnya: kembar dari $6$; kembar dari $8$; kembar dari
$10$. Lalu separuhnya: separuh dari $12$; separuh dari $20$.
\end{exercise}

\begin{exercise}[$\star$]\label{exo:g2:mult:6}
Hitunglah tabel $3$ dengan membilang loncat: $3, 6, 9, ?, ?, ?$.
Berapa $3 \times 5$?
\end{exercise}

\begin{exercise}[$\star$]\label{exo:g2:mult:7}
``Ada $4$ meja dengan $5$ kursi di tiap meja. Berapa kursi
seluruhnya?'' Tulislah perkalian dan jawabannya.
\end{exercise}

\begin{exercise}[$\star$]\label{exo:g2:mult:8}
Seekor laba-laba punya $8$ kaki. Berapa kaki $3$ ekor laba-laba?
(Bilanglah loncat $8$: $8, 16, \dots$)
\end{exercise}

\begin{exercise}[$\star$]\label{exo:g2:mult:9}
Dua cara menyusun $12$ titik dalam baris yang sama banyak: gambarlah
$2$ baris berisi $6$, lalu $3$ baris berisi $4$. Tulislah perkaliannya masing-masing.
\end{exercise}

\begin{exercise}[$\star\star$]\label{exo:g2:mult:10}
Kue dijual dalam kotak berisi $5$. Mia perlu $18$ kue untuk sebuah
pesta. Berapa kotak yang harus ia beli? (Bilanglah loncat $5$ sampai melewati $18$.)

\end{exercise}
